\documentclass[12pt]{article}
\usepackage[utf8]{inputenc}
\usepackage[T1]{fontenc}
\usepackage[italian]{babel}
\usepackage{amsmath, amssymb, amsthm}
\usepackage[margin=1in]{geometry}
\usepackage{comment}
\usepackage{graphicx}
\usepackage{verbatim}
\usepackage{tikz}
\usetikzlibrary{arrows.meta, positioning}
\usepackage{hyperref}
\usepackage{amsmath,amsfonts,amssymb}
\usepackage{geometry}
\geometry{a4paper, margin=1in}
\usepackage{graphicx}
% Pacchetti per il codice (minted + tcolorbox)
\usepackage[most]{tcolorbox}
\tcbuselibrary{minted}
\usepackage{minted}
\usepackage{xcolor}

% ------- AGGIUNTA PER SFONDO NERO E TESTO BIANCO -------
%\pagecolor{black}
%\color{white}
%\hypersetup{colorlinks=true, linkcolor=white, urlcolor=white, citecolor=white}
% -------------------------------------------------------

% ... resto del preambolo invariato ...

\newtcblisting{pythoncode}[2][]{
    listing engine=minted,
    minted language=python,
    minted options={
        style=monokai,
        fontsize=\small,
        breaklines=true,
        encoding=utf8,
        autogobble=true
    },
    colback=black!95!white,
    colframe=gray!50!black,
    listing only,
    left=5mm,
    enhanced,
    breakable,           % <--- fondamentale: permette lo spezzamento a fine pagina
    sharp corners,       % opzionale: rende gli angoli squadrati (più simile a VSCode)
    title=#2,
    #1
}

\hypersetup{colorlinks=true, linkcolor=blue, urlcolor=blue, citecolor=blue}

\newtheorem*{definizione}{Definizione}

\title{\bfseries High order FD discretization of Burgers' equation in 1D with WENO}
\author{Calcolo Scientifico (Progetto 14)}
\date{}

\begin{document}

\maketitle

\begin{abstract}
\noindent Questo documento raccoglie la teoria e le dimostrazioni alla base del progetto: la costruzione di uno schema numerico ad alta risoluzione (WENO5 + SSP-RK3) per l'equazione di Burgers non lineare, e il suo confronto con la soluzione esatta calcolata tramite il metodo delle caratteristiche e la condizione di entropia di Oleinik. Si parte dalla derivazione della condizione di salto di Rankine-Hugoniot, si mostra perché gli schemi classici basati sullo sviluppo di Taylor (Lax-Wendroff diretto) falliscono in presenza di urti, si introduce la forma conservativa come rimedio, e infine si costruisce nel dettaglio il flusso numerico WENO5.
\end{abstract}

\tableofcontents
\newpage

\section{Introduzione}

Sappiamo $u(x,0)$ per ogni $x$, e da quei valori vogliamo calcolare $u_j^{n+1}$ nel punto $x_j$ al tempo successivo. Per farlo dobbiamo usare i valori nei punti vicini $x_{j-1}$ e $x_{j+1}$ al tempo precedente, perché la derivata spaziale lega i punti tra loro: un metodo che considera più punti vicini ottiene generalmente un ordine più alto (errore minore) sulle soluzioni regolari.

Il modello che studiamo è l'equazione di Burgers non viscosa,
\[
u_t + f(u)_x = 0, \qquad f(u) = \frac{1}{2}u^2,
\]
il prototipo più semplice di legge di conservazione scalare non lineare. Nonostante la sua semplicità, l'equazione di Burgers sviluppa discontinuità (urti) in tempo finito anche a partire da dati iniziali lisci, ed è quindi il banco di prova naturale per qualunque metodo numerico pensato per leggi di conservazione più complesse (gasdinamica, traffico, ecc.).

\paragraph{Obiettivo del progetto.} Costruire uno schema alle differenze finite di alto ordine (WENO5 in spazio, SSP-RK3 in tempo) per approssimare $u(x,t)$, confrontarlo con la soluzione esatta ottenuta con il metodo delle caratteristiche e la condizione entropica di Oleinik, e verificarne numericamente l'ordine di convergenza rispetto a schemi di ordine più basso (upwind, WENO3).

\paragraph{Struttura del documento.} Le Sezioni~2--4 sviluppano la teoria necessaria: formazione degli urti, condizione di Rankine-Hugoniot, motivo per cui gli schemi basati su Taylor falliscono e necessità della forma conservativa. La Sezione~5 costruisce il flusso numerico WENO5. La Sezione~6 descrive come si calcola la soluzione esatta nel codice. \begin{comment}La Sezione~7 descrive l'implementazione e il test di convergenza\end{comment}.

\paragraph{Nota di versione.} In questa versione la formula della soluzione esatta è stata
corretta: al posto della regola euristica
\(u(x,t)=\min\{\,u_0(y): y+u_0(y)t=x\,\}\) — che con più urti che interagiscono non
seleziona la soluzione entropica — si usa la formula di \textbf{Lax--Oleinik}
(Sezione~2.4), che prende il minimo del funzionale \(J(y)=G(y)+t\,f^*((x-y)/t)\).
La correzione è motivata e verificata numericamente nelle Sezioni~2.4, 6.1, 6.4, 6.5 e
dal test \texttt{test\_soluzione\_esatta.py}; il codice di riferimento è
\texttt{codice/burgers\_weno.py} (listato in Appendice).

\newpage

\section{Formazione degli urti e soluzioni deboli}

\subsection{Il metodo delle caratteristiche}

Per l'equazione di Burgers in forma non conservativa,
\[
u_t + u u_x = 0,
\]
cerchiamo curve $x(t)$ nel piano spazio-tempo lungo le quali $u$ resta costante. Imponendo che la derivata totale di $u(x(t),t)$ sia nulla,
\[
\frac{d}{dt}u(x(t),t) = u_t + u_x \frac{dx}{dt} = 0,
\]
e confrontando con l'equazione, si vede che basta scegliere $\frac{dx}{dt} = u$. Lungo tali curve $u$ è costante, quindi anche la loro pendenza $\frac{dx}{dt}=u$ è costante ($\frac{dx}{dt}=u_0(x_0)$): le caratteristiche sono rette (basta integrare)
\[
x(t) = x_0 + u_0(x_0)\,t,
\]
dove $x_0$ è la posizione iniziale e $u_0(x_0)$ il valore iniziale della soluzione in quel punto. Il valore della soluzione in $(x,t)$ è quindi $u(x,t) = u_0(x_0)$, con $x_0$ soluzione dell'equazione implicita
\[
x_0 + u_0(x_0)\,t = x. \tag{2.1}
\]

\subsection{Tempo di rottura}

Poiché la velocità di propagazione $u_0(x_0)$ dipende dal punto di partenza, due caratteristiche uscenti da $x_1<x_2$ possono incontrarsi: ciò accade quando $u_0$ è decrescente tra $x_1$ e $x_2$. Il tempo di rottura è
\[
T_b = \min_{x\,:\,u_0'(x)<0} \left(-\frac{1}{u_0'(x)}\right),
\]
il primo istante in cui la soluzione, ancora univoca, sviluppa pendenza infinita. Dopo $T_b$ l'equazione (2.1) può avere più soluzioni $x_0$ per lo stesso $x$: la soluzione classica cessa di esistere e occorre un concetto più debole di soluzione.
Per una derivazione matematica della formula consideriamo due punti iniziali \(x_1 < x_2\). Le loro caratteristiche sono
\[
X_1(t) = x_1 + u_0(x_1) t, \qquad X_2(t) = x_2 + u_0(x_2) t.
\]
Imponiamo l'incrocio \(X_1(t) = X_2(t)\):
\[
x_1 + u_0(x_1) t = x_2 + u_0(x_2) t
\;\Longrightarrow\;
t = \frac{x_2 - x_1}{u_0(x_1) - u_0(x_2)}.
\]
Se \(u_0\) è decrescente (\(u_0(x_1) > u_0(x_2)\)), allora \(t > 0\). Per ottenere il tempo in cui due caratteristiche infinitesimamente vicine si incontrano, poniamo \(x_2 = x_1 + \Delta x\) con \(\Delta x \to 0\) e sviluppiamo \(u_0(x_2) \approx u_0(x_1) + u_0'(x_1) \Delta x\):
\[
t = \frac{\Delta x}{u_0(x_1) - [u_0(x_1) + u_0'(x_1) \Delta x]} = \frac{\Delta x}{- u_0'(x_1) \Delta x}
= -\frac{1}{u_0'(x_1)}.
\]
Il \textbf{tempo di rottura} \(T_b\) è il primo istante in cui due caratteristiche si incontrano, ovvero il minimo di questi tempi tra tutti i punti in cui \(u_0'(x) < 0\):
\[
T_b = \min_{x \,:\, u_0'(x) < 0} \left( -\frac{1}{u_0'(x)} \right).
\]

In parole povere il tempo di rottura è l'istante in cui un'onda (o un profilo di densità) diventa così ripida da "spezzarsi" formando un fronte d'urto.

Prima di quell'istante la curva è liscia e regolare; dopo inizia a formarsi una discontinuità, come quando un'onda del mare si accavalla su se stessa e si rompe.

\subsection{Soluzioni deboli}

Dopo il tempo di rottura $T_b$, le derivate $u_t, u_x$ non sono definite nel punto di salto,
quindi $u_t+f(u)_x=0$ non vale più punto per punto. Tuttavia, la legge di conservazione
in \emph{forma integrale} continua ad avere senso su ogni intervallo $[x_1,x_2]$:
\[
\frac{d}{dt}\int_{x_1}^{x_2} u(x,t)\,dx = f(u(x_1,t)) - f(u(x_2,t)). \tag{2.2}
\]

Per definire rigorosamente le soluzioni che soddisfano (2.2) anche in presenza
di discontinuità, si introduce il concetto di \textbf{soluzione debole}.
Sia \(\phi(x,t)\) una funzione test liscia a supporto compatto. Moltiplicando l'equazione per \(\phi\) e integrando su tutto lo spazio-tempo:

\[
\int_0^\infty \int_{-\infty}^{+\infty} (u_t + f(u)_x)\,\phi \, dx\,dt = 0.
\]

Integrando per parti in \(t\):

\[
\int_0^\infty \int_{-\infty}^{+\infty} u_t\,\phi \, dx\,dt 
= - \int_{-\infty}^{+\infty} u_0(x)\,\phi(x,0) \, dx - \int_0^\infty \int_{-\infty}^{+\infty} u\,\phi_t \, dx\,dt.
\]

Integrando per parti in \(x\):

\[
\int_0^\infty \int_{-\infty}^{+\infty} f(u)_x\,\phi \, dx\,dt
= - \int_0^\infty \int_{-\infty}^{+\infty} f(u)\,\phi_x \, dx\,dt.
\]

Sommando i due contributi:

\[
- \int_{-\infty}^{+\infty} u_0(x)\,\phi(x,0) \, dx
- \int_0^\infty \int_{-\infty}^{+\infty} u\,\phi_t \, dx\,dt
- \int_0^\infty \int_{-\infty}^{+\infty} f(u)\,\phi_x \, dx\,dt = 0.
\]

Moltiplicando per \(-1\):
\[
\int_0^\infty\int_{-\infty}^{+\infty} \bigl(u\phi_t + f(u)\phi_x\bigr)\,dx\,dt
+ \int_{-\infty}^{+\infty} u_0(x)\phi(x,0)\,dx = 0.
\tag{2.3}
\]

Questa formulazione è più debole perché non richiede che $u$ sia derivabile:
le derivate sono state trasferite sulla funzione test $\phi$, che è liscia.
Ogni soluzione classica (regolare) soddisfa (2.3); viceversa, una funzione
che soddisfa (2.3) per ogni $\phi$ è una \textbf{soluzione debole}.

\paragraph{Condizione di Rankine-Hugoniot.}
Per ricavare esplicitamente la condizione di salto utilizziamo la forma integrale (2.2) equivalente alla legge di conservazione integrale (2.2) su ogni intervallo \([x_1,x_2]\). Supponiamo che l'urto sia localizzato in \(x=\sigma(t)\) all'interno dell'intervallo \([x_1,x_2]\), e scomponiamo l'integrale spaziale nelle due porzioni lisce separate dalla discontinuità, applicando la regola di Leibniz per la derivazione sotto segno di integrale sugli estremi mobili. Nelle regioni in cui \(u\) è regolare vale l'equazione differenziale \(u_t=-f(u)_x\); sostituendo, i termini interni si semplificano e si ottiene la condizione algebrica (2.4):

\[
s = \frac{f(u^-)-f(u^+)}{u^- - u^+}, \qquad s=\sigma'(t). \tag{2.4}
\]

Si parte dalla legge di conservazione integrale (valida anche con urti):
\[
\frac{d}{dt} \int_{x_1}^{x_2} u(x,t) \, dx = f(u(x_1,t)) - f(u(x_2,t)). \tag{A}
\]

Sia $\Gamma: x = \sigma(t)$ la curva d'urto, con $x_1 < \sigma(t) < x_2$.
Definiamo i valori della soluzione appena a sinistra e appena a destra della curva di discontinuità:

\[
u^- := u(\sigma(t)^-, t) = \lim_{\varepsilon \to 0^+} u(\sigma(t) - \varepsilon, t), \qquad
u^+ := u(\sigma(t)^+, t) = \lim_{\varepsilon \to 0^+} u(\sigma(t) + \varepsilon, t).
\]

Dove \(\sigma(t)^-\) indica il limite sinistro della coordinata spaziale (cioè \(x \to \sigma(t)^-\)) e \(\sigma(t)^+\) il limite destro (\(x \to \sigma(t)^+\)).
Regola di Leibniz (derivata sotto integrale con estremi mobili):
\[
\frac{d}{dt} \int_{x_1}^{\sigma(t)} u \, dx = u(\sigma^-,t)\,\sigma'(t) + \int_{x_1}^{\sigma(t)} u_t \, dx,
\]
\[
\frac{d}{dt} \int_{\sigma(t)}^{x_2} u \, dx = -u(\sigma^+,t)\,\sigma'(t) + \int_{\sigma(t)}^{x_2} u_t \, dx.
\]

Sostituendo in (A), con $u^- = u(\sigma^-,t)$ e $u^+ = u(\sigma^+,t)$:
\[
(u^- - u^+) \sigma' + \int_{x_1}^{\sigma} u_t \, dx + \int_{\sigma}^{x_2} u_t \, dx = f(u(x_1)) - f(u(x_2)). \tag{B}
\]

Nelle regioni lisce ($x \neq \sigma(t)$) vale $u_t = -f(u)_x$. Quindi:
\[
\int_{x_1}^{\sigma} u_t \, dx = -\int_{x_1}^{\sigma} f(u)_x \, dx = -f(u^-) + f(u(x_1)),
\]
\[
\int_{\sigma}^{x_2} u_t \, dx = -\int_{\sigma}^{x_2} f(u)_x \, dx = -f(u(x_2)) + f(u^+).
\]

Sostituendo in (B):
\[
(u^- - u^+) \sigma' - f(u^-) + f(u(x_1)) - f(u(x_2)) + f(u^+) = f(u(x_1)) - f(u(x_2)).
\]

Semplificando i termini $f(u(x_1)) - f(u(x_2))$ da entrambi i lati:
\[
(u^- - u^+) \sigma' - f(u^-) + f(u^+) = 0.
\]

Da cui si ricava la condizione di Rankine--Hugoniot:
\[
{\displaystyle \sigma'(t) = \frac{f(u^-) - f(u^+)}{u^- - u^+}} \quad \text{oppure} \quad s = \frac{\Delta f}{\Delta u}.
\]
Per l'equazione di Burgers, $f(u)=\frac12 u^2$, quindi
\[
s = \frac{\frac12(u^-)^2-\frac12(u^+)^2}{u^- - u^+}
  = \frac{u^-+u^+}{2}.
\]

\subsection{Condizioni di entropia}

Per $u^->u^+$ la soluzione debole corrisponde a un urto fisico; per $u^-<u^+$
esiste anche un ``urto di espansione'' non fisico. Per selezionare la soluzione
corretta si impone la \textbf{condizione di Lax}:
\[
f'(u^-) > s > f'(u^+),
\]
che per Burgers equivale a $u^->u^+$.
\newline
\textit{L'urto deve saltare da un valore alto a sinistra a uno basso a destra}
\newline
Una formulazione equivalente è la \textbf{condizione di Oleinik}: per ogni $x_2>x_1$ e $t>0$
\[
u(x_2,t)-u(x_1,t) \le \frac{x_2-x_1}{t},
\]
cioè $u(\cdot,t)$ non può crescere più rapidamente di $1/t$: è la forma integrata del
requisito che escluda gli urti di rarefazione ($u^-<u^+$).
Per flusso convesso ($f''>0$, come per Burgers) la soluzione entropica si scrive in forma
chiusa con la \textbf{formula di Lax--Oleinik} (o \emph{Hopf--Lax--Oleinik}):
\[
u(x,t) = (f^*)'\!\left(\frac{x-y^*}{t}\right), \qquad
y^* = \arg\min_{y\in\mathbb{R}}\left[\, G(y) + t\,f^*\!\left(\frac{x-y}{t}\right)\right],
\qquad G(y)=\int_{y_0}^{y} u_0(s)\,ds,
\]
dove $f^*(p)=\sup_u\,[\,pu-f(u)\,]$ è la trasformata di Legendre--Fenchel di $f$. Per
Burgers $f^*(p)=p^2/2$, quindi $(f^*)'(p)=p$ e la formula diventa
\[
u(x,t) = \frac{x-y^*}{t}, \qquad
y^* = \arg\min_{y\in\mathbb{R}}\left[\, G(y) + \frac{(x-y)^2}{2t}\right].
\]

\textbf{Attenzione (questo è il punto delicato, corretto in questa versione del documento).}
I punti stazionari del funzionale $J(y)=G(y)+t\,f^*((x-y)/t)$ sono esattamente le radici di
$y+u_0(y)t=x$, cioè le caratteristiche che arrivano in $(x,t)$: restringersi a quelle radici
è corretto. Quello che \emph{non} è corretto è scegliere fra i candidati quello con $u_0(y)$
minimo, come si fa talvolta in modo euristico: la selezione entropica consiste nel prendere
il candidato che \textbf{minimizza il funzionale $J$}. I due criteri coincidono quando la
soluzione è a un solo valore o quando c'è un solo urto isolato, ma \emph{divergono} appena
più urti interagiscono: in quel caso $\min\{u_0(y): y+u_0(y)t=x\}$ può pescare una
caratteristica molto lontana (con $u_0$ piccolo) e restituire un valore che non è la
soluzione fisica, mentre $J$ confronta le aree $\int u_0$ pesate dal termine di
Legendre--Fenchel.
Si veda la Sezione~6 per la verifica numerica e il test automatico allegato al codice.


\newpage

\section{L'approccio di Taylor e il suo fallimento sugli urti}

\subsection{Il problema lineare: l'equazione del trasporto}

Consideriamo l'equazione del trasporto lineare, con velocità costante $a\in\mathbb{R}$: $u_t = -au_x$. Sviluppando in serie di Taylor rispetto al tempo fino al secondo ordine,
\[
u_j^{n+1} = u_j^n + \Delta t\,(u_t)_j^n + \frac{\Delta t^2}{2}(u_{tt})_j^n + O(\Delta t^3).
\]
Il primo termine è $(u_t)_j^n=-a(u_x)_j^n$. Per il secondo, deriviamo $u_t$ rispetto al tempo e usiamo il teorema di Schwarz per scambiare $u_{xt}=(u_t)_x$:
\[
u_{tt} = \frac{\partial}{\partial t}(-au_x) = -a(u_t)_x = -a(-au_x)_x = a^2 u_{xx}.
\]
Sostituendo, si ottiene la forma differenziale continua del secondo ordine:
\[
u_j^{n+1} = u_j^n - \Delta t\,a(u_x)_j^n + \frac{\Delta t^2}{2}a^2(u_{xx})_j^n.
\]
Discretizzando le derivate spaziali con differenze centrate,
\[
(u_x)_j^n \approx \frac{u_{j+1}^n-u_{j-1}^n}{2\Delta x}, \qquad
(u_{xx})_j^n \approx \frac{u_{j+1}^n-2u_j^n+u_{j-1}^n}{\Delta x^2},
\]
si ottiene lo schema esplicito di Lax-Wendroff per il singolo nodo $j$:
\[
u_j^{n+1} = u_j^n - \frac{a\Delta t}{2\Delta x}\bigl(u_{j+1}^n-u_{j-1}^n\bigr) + \frac{a^2\Delta t^2}{2\Delta x^2}\bigl(u_{j+1}^n-2u_j^n+u_{j-1}^n\bigr).
\]
Questo schema è l'estensione naturale del metodo di Taylor al secondo ordine per il caso lineare, ed è esatto fino a $O(\Delta t^3)$ perché la soluzione $u(x,t)=u_0(x-at)$ resta liscia quanto lo è il dato iniziale.

\newpage
\subsection{Il problema non lineare: Burgers (approccio diretto)}

Per Burgers, in forma non conservativa $u_t=-uu_x$. Applichiamo lo stesso procedimento di Taylor. Il termine del primo ordine è $(u_t)_j^n=-u_j^n(u_x)_j^n$. Per il secondo ordine deriviamo ancora rispetto al tempo:
\[
u_{tt} = \frac{\partial}{\partial t}(-uu_x) = -u_tu_x-uu_{xt}.
\]
Sfruttando Schwarz ($u_{xt}=(u_t)_x$) e sostituendo $u_t=-uu_x$ e $(u_t)_x=(-uu_x)_x=-(u_x)^2-uu_{xx}$, si ottiene:
\[
u_{tt} = -(-uu_x)u_x - u\left(-(u_x)^2 - uu_{xx}\right)
      = u(u_x)^2 + u(u_x)^2 + u^2u_{xx}
      = 2u(u_x)^2 + u^2u_{xx}.
\]
Sostituendo $u_t$ e $u_{tt}$ nello sviluppo di Taylor:
\[
u_j^{n+1} = u_j^n - \Delta t\bigl(u_j^n u_x\bigr) + \frac{\Delta t^2}{2}\left(2u_j^n (u_x)^2 + (u_j^n)^2 u_{xx}\right).
\]
Discretizzando $(u_x)_j^n \approx \frac{u_{j+1}^n-u_{j-1}^n}{2\Delta x}$ e $(u_{xx})_j^n \approx \frac{u_{j+1}^n-2u_j^n+u_{j-1}^n}{\Delta x^2}$, si ottiene lo schema esplicito
\[
\begin{aligned}
u_j^{n+1} = u_j^n &- \frac{\Delta t}{2\Delta x}u_j^n\bigl(u_{j+1}^n-u_{j-1}^n\bigr) \\
&+ \frac{\Delta t^2}{2}\left[
2u_j^n \left(\frac{u_{j+1}^n-u_{j-1}^n}{2\Delta x}\right)^2
+ (u_j^n)^2 \frac{u_{j+1}^n-2u_j^n+u_{j-1}^n}{\Delta x^2}
\right].
\end{aligned}
\]
Questo schema è la traduzione letterale del ``trucco di Lax-Wendroff'' al caso non lineare: approssima direttamente il valore puntuale $u_j^{n+1}$ tramite Taylor, senza passare per la forma integrale.

\subsection{Fallimento dell'approccio diretto di Taylor sugli urti}

\textit{Lo sviluppo di Taylor richiede la regolarità di $u$, ma in un urto $u_x$ diverge e le derivate superiori non esistono. Lo schema genera oscillazioni di Gibbs, perde consistenza (l'errore di troncamento contiene $u_{xxx}$, infinito all'urto) e viola i vincoli fisici. Il problema più grave è la non conservazione della massa discreta, che produce una velocità dell'urto $s_{\text{num}} \neq s$ errata, anche al raffinare la griglia. Il rimedio è sostituire l'approssimazione puntuale con quella della forma integrale, valida anche attraverso le discontinuità.}
\newpage

\section{Forma conservativa}

\subsection{Derivazione dalla forma integrale}

Prendiamo una cella spaziale $[x_{j-1/2},x_{j+1/2}]$ e integriamo $u_t+f(u)_x=0$ su questa cella e sull'intervallo temporale $[t_n,t_{n+1}]$:
\[
\int_{t_n}^{t_{n+1}}\int_{x_{j-1/2}}^{x_{j+1/2}} \bigl(u_t+f(u)_x\bigr)\,dx\,dt = 0.
\]
Scambiando l'ordine di integrazione:
\[
\int_{x_{j-1/2}}^{x_{j+1/2}}\Bigl(\int_{t_n}^{t_{n+1}} u_t\,dt\Bigr)dx + \int_{t_n}^{t_{n+1}}\Bigl(\int_{x_{j-1/2}}^{x_{j+1/2}} f(u)_x\,dx\Bigr)dt = 0.
\]
Il primo integrale in $t$ è $u(x,t_{n+1})-u(x,t_n)$; il secondo in $x$ è $f(u(x_{j+1/2},t))-f(u(x_{j-1/2},t))$ per il teorema fondamentale del calcolo. Dividendo per $\Delta x$ e definendo la media di cella
\[
U_j^n = \frac{1}{\Delta x}\int_{x_{j-1/2}}^{x_{j+1/2}} u(x,t_n)\,dx,
\]
si ottiene
\[
U_j^{n+1} = U_j^n - \frac{1}{\Delta x}\left(\int_{t_n}^{t_{n+1}} f(u(x_{j+1/2},t))\,dt - \int_{t_n}^{t_{n+1}} f(u(x_{j-1/2},t))\,dt\right).
\]
La quantità $\frac{1}{\Delta t}\int_{t_n}^{t_{n+1}} f(u(x_{j+1/2},t))\,dt$ è il \textbf{flusso medio} attraverso il bordo destro della cella; la chiamiamo $F_{j+1/2}^n$. Con la regola del punto medio (accurata al second'ordine),
\[
F_{j+1/2}^n \approx f\bigl(u(x_{j+1/2},t_{n+1/2})\bigr) + O(\Delta t^2).
\]

\newpage
\subsection{Il flusso di Lax-Wendroff in forma conservativa}

Sviluppiamo in Taylor \emph{in tempo} la funzione $g(x,t)=f(u(x,t))$ attorno a $t_n$, fermandoci al prim'ordine (l'errore sarà $O(\Delta t^2)$ dopo aver moltiplicato per $\Delta t$):
\[
g(x_{j+1/2},t_{n+1/2}) = g(x_{j+1/2},t_n) + \frac{\Delta t}{2}g_t(x_{j+1/2},t_n) + O(\Delta t^2).
\]
Per la regola della catena $g_t=f'(u)u_t$; usando $u_t=-f(u)_x=-g_x$ e ponendo $a(u)=f'(u)$, si ha $g_t=-a(u)g_x$. Quindi
\[
F_{j+1/2}^n = g(x_{j+1/2},t_n) - \frac{\Delta t}{2}a(u)\,g_x(x_{j+1/2},t_n) + O(\Delta t^2).
\]
Approssimando i termini al bordo con i valori nelle celle $j$ e $j+1$: la media aritmetica per $g$, il rapporto incrementale centrato per $g_x$, e la media dei due stati per $a(u)$,
\[
g(x_{j+1/2},t_n)\approx\frac{g_j+g_{j+1}}{2}, \qquad
g_x(x_{j+1/2},t_n)\approx\frac{g_{j+1}-g_j}{\Delta x}, \qquad
a_{j+1/2}=\frac{U_j^n+U_{j+1}^n}{2},
\]
si arriva alla \textbf{formula del flusso di Lax-Wendroff in forma conservativa}:
\[
F_{j+1/2}^n = \frac{g_j+g_{j+1}}{2} - \frac{\Delta t}{2\Delta x}\,a_{j+1/2}\,(g_{j+1}-g_j).
\]
La differenza sostanziale rispetto all'approccio diretto della Sezione~3 è che qui si approssima il \emph{flusso integrato al bordo} $F_{j+1/2}^n=f(u(x_{j+1/2},t_{n+1/2}))$, quantità che esiste sempre perché deriva dalla forma integrale, mentre l'approccio diretto approssima la soluzione puntuale $u_j^{n+1}$, che sull'urto non esiste in senso classico.

\newpage

\subsection{Conservazione della massa: approccio conservativo}

Consideriamo lo schema $u_j^{n+1}=u_j^n-\frac{\Delta t}{\Delta x}(F_{j+1/2}^n-F_{j-1/2}^n)$ su un dominio periodico con $J$ celle, e definiamo la massa discreta $\mathcal{M}^n=\Delta x\sum_{j=1}^J u_j^n$. Allora
\[
\mathcal{M}^{n+1} = \Delta x\sum_{j=1}^J u_j^n - \Delta t\sum_{j=1}^J\bigl(F_{j+1/2}^n-F_{j-1/2}^n\bigr) = \mathcal{M}^n - \Delta t\sum_{j=1}^J\bigl(F_{j+1/2}^n-F_{j-1/2}^n\bigr).
\]
La somma dei flussi è telescopica:
\[
\sum_{j=1}^J\bigl(F_{j+1/2}^n-F_{j-1/2}^n\bigr) = F_{J+1/2}^n - F_{1-1/2}^n,
\]
e per periodicità $F_{J+1/2}^n=F_{1-1/2}^n$, quindi la somma è nulla e $\mathcal{M}^{n+1}=\mathcal{M}^n$. Poiché ciò vale per ogni $n$, $\mathcal{M}^n=\mathcal{M}^0$ per ogni $n$: la forma conservativa preserva \emph{esattamente} il bilancio di massa discreto, indipendentemente da quanto sia ripida la soluzione. Per il teorema di Lax-Wendroff (Sezione~4.5), se uno schema conservativo e consistente converge, converge a una soluzione debole corretta: la velocità dell'urto numerico coincide con quella di Rankine-Hugoniot.

\newpage\subsection{Non conservazione della massa: approccio diretto}

Consideriamo ora lo schema ottenuto in Sezione 3.2 dallo sviluppo di Taylor diretto sulla forma non conservativa, discretizzando le derivate spaziali con differenze centrate, si ottiene:

\[
u_j^{n+1} = u_j^n - \frac{\Delta t}{2\Delta x} u_j (u_{j+1} - u_{j-1}) 
+ \frac{\Delta t^2}{2\Delta x^2} u_j^2 (u_{j+1} - 2u_j + u_{j-1}) 
+ \frac{\Delta t^2}{4\Delta x^2} u_j (u_{j+1} - u_{j-1})^2.
\]

Calcoliamo ora la massa discreta \(\mathcal{M}^n = \Delta x \sum_{j=1}^J u_j^n\). La variazione di massa dal passo \(n\) al passo \(n+1\) è:

\[
\mathcal{M}^{n+1} - \mathcal{M}^n = 
- \frac{\Delta t}{2} \sum_{j=1}^J u_j (u_{j+1} - u_{j-1}) 
+ \frac{\Delta t^2}{2\Delta x} \sum_{j=1}^J u_j^2 (u_{j+1} - 2u_j + u_{j-1}) 
+ \frac{\Delta t^2}{4\Delta x} \sum_{j=1}^J u_j (u_{j+1} - u_{j-1})^2.
\]

Per periodicità, il termine \(O(\Delta t)\) si annulla (è una traslazione di indici):
\[
\sum_{j=1}^J u_j^nu_{j+1}^n - \sum_{j=1}^J u_j^nu_{j-1}^n = 0.
\]

Per i termini \(O(\Delta t^2)\), definiamo la quantità
\[
\mathcal{R}(u^n) = 
\sum_{j=1}^J u_j^2 (u_{j+1} - 2u_j + u_{j-1}) 
+ \frac{1}{2} \sum_{j=1}^J u_j (u_{j+1} - u_{j-1})^2.
\]

Espandendo il primo termine e traslando gli indici (\(\sum_j u_j^2 u_{j-1} = \sum_j u_{j+1}^2 u_j\)), si ha:
\[
\sum_{j=1}^J u_j^2 (u_{j+1} - 2u_j + u_{j-1}) 
= \sum_{j=1}^J \Bigl[u_j u_{j+1}(u_j+u_{j+1}) - 2u_j^3\Bigr].
\]

Il secondo termine non può essere riscritto come una differenza telescopica e non si semplifica con il primo in modo da dare zero identicamente. In generale,
\[
\mathcal{R}(u^n) \neq 0,
\]
come si verifica immediatamente anche su profili semplici non costanti (ad esempio, su una griglia periodica con valori \(u=[1,2,3]\), \(\mathcal{R}\neq 0\)). 

Pertanto, la variazione di massa è
\[
\mathcal{M}^{n+1} = \mathcal{M}^n + \frac{\Delta t^2}{2\Delta x}\mathcal{R}(u^n) \neq \mathcal{M}^n.
\]

In presenza di un urto, le differenze \(u_{j+1}^n-u_j^n\) sono \(O(1)\) indipendentemente dal raffinamento della griglia, quindi \(\mathcal{R}(u^n)=O(1)\) non tende a zero. Sotto la condizione CFL (\(\lambda=\Delta t/\Delta x\) costante), l'errore di massa per singolo passo temporale scala come
\[
\frac{\Delta t^2}{2\Delta x}\mathcal{R}(u^n) = \frac{\Delta t}{2}\lambda\,\mathcal{R}(u^n) = O(\Delta t).
\]
Poiché per raggiungere un tempo finale \(T\) sono necessari \(N_t = O(1/\Delta t)\) passi temporali, l'errore di massa accumulato è \(O(1)\): si tratta di un deficit permanente, che non scompare al raffinare la griglia. Questo viola la condizione di Rankine-Hugoniot a livello discreto e costringe l'urto numerico a propagarsi a una velocità \(s_{\text{num}}\neq s\) errata, anche nel limite di griglia fine.
\newpage

\subsection{Generalizzazione: il caso lineare come caso particolare}

Verifichiamo che la formula del flusso conservativo generalizza correttamente il metodo di Lax-Wendroff lineare della Sezione~3.1. Per $f(u)=au$ con $a$ costante si ha $g_j=au_j$ e $a(u)=f'(u)=a$. Sostituendo nella formula del flusso,
\[
F_{j+1/2}^n = \frac{g_j+g_{j+1}}{2} - \frac{\Delta t}{2\Delta x}a_{j+1/2}(g_{j+1}-g_j),
\]
si ottiene
\[
F_{j+1/2} = \frac{a}{2}(u_j+u_{j+1}) - \frac{a^2\Delta t}{2\Delta x}(u_{j+1}-u_j).
\]
Calcolando $F_{j+1/2}-F_{j-1/2}$ e sostituendo in $U_j^{n+1}=U_j^n-\frac{\Delta t}{\Delta x}(F_{j+1/2}-F_{j-1/2})$ si ritrova esattamente lo schema di Lax-Wendroff lineare:
\[
U_j^{n+1} = U_j^n - \frac{a\Delta t}{2\Delta x}(U_{j+1}-U_{j-1}) + \frac{a^2\Delta t^2}{2\Delta x^2}(U_{j+1}-2U_j+U_{j-1}).
\]
Per $f(u)$ lineare, dunque, l'approccio diretto e quello conservativo coincidono (grazie alla linearità): è per questo che nella Sezione~3.1 il metodo di Taylor diretto è ineccepibile. Solo quando $f(u)$ è non lineare le due strade divergono, e solo la forma conservativa garantisce la corretta velocità dell'urto.

\subsection{Perché la forma conservativa basta: cancellazione telescopica}

Il punto cruciale, che riassume l'intera sezione, è puramente algebrico. Scrivendo l'aggiornamento come $U_j^{n+1}=U_j^n-\frac{\Delta t}{\Delta x}(F_{j+1/2}-F_{j-1/2})$ e sommando su tutte le celle (bordo periodico o flussi nulli), tutti i flussi interni si cancellano a coppie:
\[
\sum_j U_j^{n+1}\Delta x = \sum_j U_j^n\Delta x.
\]
La quantità totale si conserva \emph{esattamente}, qualunque sia la ripidità della soluzione: è questa proprietà, e non l'accuratezza puntuale, a garantire la corretta velocità dell'urto secondo Rankine-Hugoniot. Questo principio -- forma conservativa $\Rightarrow$ massa esatta $\Rightarrow$ urto corretto -- è lo stesso che useremo per costruire il metodo WENO nella prossima sezione: cambierà soltanto \emph{come} si approssima il flusso $F_{j+1/2}$, non la struttura conservativa dello schema.

\newpage

\section{Il metodo WENO (Weighted Essentially Non-Oscillatory)}

\subsection{Motivazione}

Gli schemi conservativi del secondo ordine come Lax-Wendroff (Sezione~4) garantiscono la corretta velocità dell'urto, ma restano sviluppi di Taylor \emph{applicati al flusso}: in prossimità di una discontinuità producono comunque oscillazioni di Gibbs, perché il polinomio ricostruito tenta di interpolare punti da entrambi i lati del salto. Gli schemi del primo ordine come Godunov sono monotoni (non oscillano) ma eccessivamente dissipativi: l'urto viene spalmato su molte celle.

L'idea dei metodi ENO (Essentially Non-Oscillatory) è: se uno stencil attraversa una discontinuità, non usarlo per la ricostruzione. Per ogni bordo di cella si valutano più stencil candidati e si sceglie quello più liscio. I metodi WENO generalizzano l'idea usando una \emph{media pesata} di tutti gli stencil, privilegiando quelli lisci: nelle zone regolari si recupera così l'ordine massimo possibile, mentre vicino agli urti i pesi degli stencil ``cattivi'' si annullano automaticamente.

\subsection{Struttura del metodo: separazione tra spazio e tempo}
Il metodo si articola in due fasi indipendenti.

\textbf{Fase 1 (WENO, spazio):} si sostituisce \(f(u)_x\) con un operatore numerico \(L(U)_j\), ottenendo per ogni cella un'equazione differenziale ordinaria nel tempo,
\[
\frac{du_j}{dt} = L(U(t))_j = \underbrace{-\frac{1}{\Delta x}\bigl(\hat f_{j+1/2}(t)-\hat f_{j-1/2}(t)\bigr)}_{\textit{Stima numerica di } -f(u)_x \textit{ in } x=x_j}\]
dove \(U(t)=(u_1(t),\dots,u_N(t))^T\) è il vettore delle incognite all'istante \(t\). 
Si noti che \(L\) \textbf{non dipende esplicitamente} da \(t\) (nella formula non compare mai la variabile \(t\)); la dipendenza dal tempo è \textbf{solo implicita}, perché \(L\) viene valutata sui valori \(U(t)\) che cambiano nel tempo.
Il problema diventa quindi un sistema di \(N\) ODE accoppiate (problema di Cauchy):
\[
\frac{dU}{dt} = L(U(t)), \qquad U(0)=U_0,
\]

\textbf{Fase 2 (Runge-Kutta, tempo):} l'integrale esatto 
\(U(t_{n+1})=U(t_n)+\int_{t_n}^{t_{n+1}} L(U(\tau))\,d\tau\) 
va approssimato con una quadratura temporale di ordine coerente con quello spaziale. Eulero esplicito darebbe solo ordine 1, sprecando l'accuratezza di WENO5; si usa quindi il metodo \textbf{SSP-RK3} (Strong Stability Preserving Runge-Kutta a 3 stadi). 
Ponendo \(U^n = U(t_n)\), e indicando con \(U^{(1)} \approx U(t_n+\Delta t)\) e \(U^{(2)} \approx U(t_n+\tfrac12\Delta t)\) le soluzioni approssimate ai tempi intermedi, gli stadi sono:
\[
\begin{aligned}
U^{(1)} &= U^n + \Delta t\,L(U^n), \\[4pt]
U^{(2)} &= \tfrac34 U^n + \tfrac14 (U^{(1)} + \Delta t\,L(U^{(1)})), \\[4pt]
U^{n+1} &= \tfrac13 U^n + \tfrac23 (U^{(2)} + \Delta t\,L(U^{(2)})).
\end{aligned}
\]
Poiché \(\Delta t \propto \Delta x\), l'errore temporale scala come \(\Delta x^3\), mentre quello spaziale come \(\Delta x^5\). Per \(\Delta x\) sufficientemente piccolo, il termine \(\Delta x^3\) domina, quindi l'ordine globale del metodo è 3.
\subsection{Flux splitting di Lax-Friedrichs}

Al tempo \(t_n\), noto il vettore \(U^n = U(t_n)\), si calcolano i flussi fisici puntuali 
\(f_j(t_n)=\frac12\bigl(U_j(t_n)\bigr)^2\). 
Si applica lo split di Lax-Friedrichs per separare l'informazione che viaggia verso destra da quella che viaggia verso sinistra:
\[
f^+(u) = \frac12\bigl(f(u)+\alpha u\bigr), \qquad 
f^-(u) = \frac12\bigl(f(u)-\alpha u\bigr), \qquad 
\alpha=\max_j|U_j(t_n)|,
\]
\textit{Questa è la scelta standard per $f^+$ ed $f^-$: garantisce che le derivate abbiano segno costante su tutto il dominio e che la loro somma restituisca esattamente il flusso originale $f(u)$. Prendiamo $\alpha = \max_j |U_j(t_n)|$ in modo che, per ogni punto della griglia, valga $-\alpha \le u_j \le \alpha$; di conseguenza, siccome $f'(u)=u$ per Burgers}
\[
(f^+)'(u_j) = \frac{1}{2}(u_j+\alpha) \ge 0, \qquad (f^-)'(u_j) = \frac{1}{2}(u_j-\alpha) \le 0 \quad \text{per ogni } j,
\]
\textit{cioè le due componenti viaggiano effettivamente in direzioni opposte e ben definite. 
Se avessimo scelto $\alpha < \max_j |u_j|$, esisterebbe almeno un punto $j$ in cui $u_j > \alpha$ (se il massimo è positivo) o $u_j < -\alpha$ (se il minimo è negativo); in quel punto avremmo $(f^-)'(u_j) > 0$ oppure $(f^+)'(u_j) < 0$, violando le condizioni di upwind e rendendo instabile la ricostruzione di WENO su quel bordo.}

\textit{Si sceglie $\alpha = \max_j |U_j|$ (il minimo valore che garantisce la proprietà per tutti i punti; valori più grandi sono ammissibili ma introducono maggiore dissipazione numerica).}
Grazie a questa proprietà, \(f^+\) si ricostruisce con uno stencil sbilanciato a sinistra (upwind per l'onda che va a destra) e \(f^-\) con uno stencil sbilanciato a destra. 
Il flusso numerico finale al bordo \(j+\tfrac12\) all'istante \(t_n\) è quindi:
\[
\hat f_{j+1/2}(t_n)=\hat f^+_{j+1/2}(t_n)+\hat f^-_{j+1/2}(t_n).
\]


Nel seguito descriviamo la ricostruzione di \(\hat f^+_{j+1/2}(t_n)\); quella di \(\hat f^-_{j+1/2}(t_n)\) è analoga con indici specchiati. 
Si ricordi che durante i sottostadi di Runge-Kutta, lo stesso procedimento verrà ripetuto valutando \(L\) sui vettori intermedi \(U^{(1)}\) e \(U^{(2)}\), producendo quindi le stime \(\hat f_{j+1/2}(t_n+\Delta t)\) e \(\hat f_{j+1/2}(t_n+\tfrac12\Delta t)\) necessarie all'integrazione temporale.

\newpage

\subsection{I tre sotto-stencil e la derivazione di $p_0$}

Per WENO5 i tre sotto-stencil per il bordo $x_{j+1/2}$ sono
\[
S_0=\{j-2,j-1,j\}, \qquad S_1=\{j-1,j,j+1\}, \qquad S_2=\{j,j+1,j+2\},
\]
e su ciascuno si costruisce un polinomio $p_r(x)$ di grado 2, valutato poi al bordo $x_{j+1/2}$. Mostriamo la derivazione di $p_0$ sullo stencil $S_0$. Nel metodo dei volumi finiti $f_j$ è una media di cella, $f_j=\frac{1}{\Delta x}\int_{x_{j-1/2}}^{x_{j+1/2}} f(x)\,dx$; cerchiamo $p(x)=ax^2+bx+c$ la cui media sulle tre celle di $S_0$ riproduca $f_{j-2},f_{j-1},f_j$. Con $\Delta x=1$, $j=0$:
\[
\int_{-5/2}^{-3/2}p\,dx=f_{-2}, \qquad \int_{-3/2}^{-1/2}p\,dx=f_{-1}, \qquad \int_{-1/2}^{1/2}p\,dx=f_0.
\]
Calcolando gli integrali si ottiene il sistema
\[
\begin{cases}
\frac{49}{12}a-2b+c=f_{-2}, \\
\frac{13}{12}a-b+c=f_{-1}, \\
\frac{1}{12}a+c=f_0.
\end{cases}
\]
Sottraendo la seconda equazione dalla prima e la terza dalla seconda:
\[
3a-b=f_{-2}-f_{-1}, \qquad a-b=f_{-1}-f_0.
\]
Sottraendo queste due tra loro: $2a=f_{-2}-2f_{-1}+f_0$, quindi
\[
a=\frac{f_{-2}-2f_{-1}+f_0}{2}, \qquad b=a-(f_{-1}-f_0)=\frac{f_{-2}-4f_{-1}+3f_0}{2}, \qquad c=f_0-\frac{a}{12}.
\]
Valutando in $x=1/2$ (il bordo $x_{j+1/2}$):
\[
p_0\left(\tfrac12\right) = \frac{a}{4}+\frac{b}{2}+c = \frac13 f_{j-2}-\frac76 f_{j-1}+\frac{11}{6}f_j.
\]
Con lo stesso procedimento, traslando gli indici sugli stencil $S_1$ e $S_2$, si ottengono gli altri due polinomi:
\[
p_1 = -\frac16 f_{j-1}+\frac56 f_j+\frac13 f_{j+1}, \qquad
p_2 = \frac13 f_j+\frac56 f_{j+1}-\frac16 f_{j+2}.
\]
Se combinati con pesi costanti $d_0=\frac{1}{10},\,d_1=\frac35,\,d_2=\frac{3}{10}$ (i \textbf{pesi lineari ottimali}, con $d_0+d_1+d_2=1$), i tre polinomi riproducono esattamente il polinomio di quarto grado che interpola tutti e cinque i punti dello stencil globale, dando ordine 5 nelle zone lisce. Vicino a un urto, però, questi pesi userebbero anche stencil che attraversano il salto: vanno resi adattativi.
Questi pesi $d_0,d_1,d_2$ si calcolano risolvendo il sistema $P_{j+1/2}=d_0p_0+d_1p_1+d_2p_2$ dove P è il polinomio che risolve \[
\begin{cases}
\int_{-5/2}^{-3/2} (\alpha x^4 + \beta x^3 + \gamma x^2 + \delta x + \eta)\,dx = f_{-2} \\
\int_{-3/2}^{-1/2} (\alpha x^4 + \beta x^3 + \gamma x^2 + \delta x + \eta)\,dx = f_{-1} \\
\int_{-1/2}^{1/2} (\alpha x^4 + \beta x^3 + \gamma x^2 + \delta x + \eta)\,dx = f_0 \\
\int_{1/2}^{3/2} (\alpha x^4 + \beta x^3 + \gamma x^2 + \delta x + \eta)\,dx = f_1 \\
\int_{3/2}^{5/2} (\alpha x^4 + \beta x^3 + \gamma x^2 + \delta x + \eta)\,dx = f_2
\end{cases}
\]

\newpage

\subsection{Indicatori di lisciatura e pesi non lineari}

Per ogni stencil $S_r$ si definisce un \textbf{indicatore di lisciatura} $\beta_r$, piccolo nelle zone regolari e grande se lo stencil attraversa una discontinuità:
\[
\begin{aligned}
\beta_0 &= \tfrac{13}{12}(f_{j-2}-2f_{j-1}+f_j)^2 + \tfrac14(f_{j-2}-4f_{j-1}+3f_j)^2, \\
\beta_1 &= \tfrac{13}{12}(f_{j-1}-2f_j+f_{j+1})^2 + \tfrac14(f_{j-1}-f_{j+1})^2, \\
\beta_2 &= \tfrac{13}{12}(f_j-2f_{j+1}+f_{j+2})^2 + \tfrac14(3f_j-4f_{j+1}+f_{j+2})^2,
\end{aligned}
\]
dove il primo addendo misura la curvatura (derivata seconda discreta) e il secondo la pendenza (derivata prima discreta) all'interno dello stencil. Mostriamo sotto la derivazione di $\beta_0$:
\newline

\[
p_0(x)=ax^2+bx+c,\qquad 
a=\frac{f_{-2}-2f_{-1}+f_0}{2},\quad
b=\frac{f_{-2}-4f_{-1}+3f_0}{2}
\]

\[
p_0'(x)=2ax+b,\qquad p_0''(x)=2a
\]

\[
\beta_0 = \underbrace{\int_{-1/2}^{1/2} (p_0')^2 dx}_{\displaystyle \frac{a^2}{3}+b^2}
      + \underbrace{\int_{-1/2}^{1/2} (p_0'')^2 dx}_{\displaystyle 4a^2}
\]

\[
\frac{a^2}{3}+b^2+4a^2
= \frac{a^2}{3}+4a^2+b^2
= \frac{13a^2}{3}+b^2
\]

\[
\frac{13a^2}{3} = \frac{13}{3}\left(\frac{f_{-2}-2f_{-1}+f_0}{2}\right)^2
= \frac{13}{12}(f_{-2}-2f_{-1}+f_0)^2
\]

\[
b^2 = \frac{1}{4}(f_{-2}-4f_{-1}+3f_0)^2
\]

\[
\beta_0 = \frac{13}{12}(f_{-2}-2f_{-1}+f_0)^2 + \frac{1}{4}(f_{-2}-4f_{-1}+3f_0)^2
\]



I pesi non lineari si costruiscono da $d_r$ e $\beta_r$:
\[
\alpha_r = \frac{d_r}{(\beta_r+\varepsilon)^2}, \qquad \omega_r=\frac{\alpha_r}{\alpha_0+\alpha_1+\alpha_2}, \qquad \varepsilon=10^{-6}\ (\text{evita la divisione per zero}).
\]
Se $\beta_r$ è piccolo (zona liscia), $\omega_r\approx d_r$ e si recupera il quinto ordine; se $\beta_r$ è grande (stencil che attraversa l'urto), $\omega_r\approx0$ e quello stencil viene escluso dalla ricostruzione. I pesi non si annullano mai esattamente (da qui `Essentially' di WENO), ma diventano trascurabili. Il flusso ricostruito al bordo è infine
\[
\hat f^+_{j+1/2} = \omega_0p_0+\omega_1p_1+\omega_2p_2.
\]
\subsection{Proprietà del metodo}

\textit{Il metodo WENO5 gode di quattro proprietà fondamentali. La conservatività è garantita dalla struttura stessa dello schema, $U_j^{n+1}=U_j^n-\frac{\Delta t}{\Delta x}(\hat f_{j+1/2}-\hat f_{j-1/2})$, per cui la cancellazione telescopica dei flussi ai bordi interni assicura la conservazione esatta della massa totale su dominio periodico (Sezione~4.6). L'assenza di oscillazioni di Gibbs deriva dal comportamento dei pesi non lineari: in prossimità di un urto, gli stencil che lo attraversano ricevono $\omega_r \approx 0$, escludendo di fatto dalla ricostruzione i dati provenienti dal lato opposto della discontinuità. Nelle zone lisce, dove gli indicatori $\beta_r$ sono confrontabili, i pesi non lineari tendono ai pesi ottimali $\omega_r \approx d_r$, e la combinazione dei tre polinomi recupera l'accuratezza di quinto ordine. Infine, la selezione della soluzione entropica è garantita dalla dissipazione numerica introdotta dallo split di Lax-Friedrichs, che orienta la ricostruzione nella direzione delle caratteristiche e seleziona la soluzione debole entropica coerente con la condizione di Oleinik (Sezione~2.4).}
\newpage


\subsection{Pseudocodice}

Lo schema completo si articola nei seguenti passaggi, ripetuti per ogni passo temporale \(t_n \to t_{n+1}\) con l'integratore SSP-RK3.

\begin{enumerate}
    \item \textbf{Calcolo dell'operatore spaziale} \(L(U)\).
    
    Dato un vettore \(U=(U_1,\dots,U_N)\) rappresentante la soluzione in tutte le celle in un dato istante, si esegue:
    
    \begin{enumerate}
        \item \textbf{Flussi fisici:} per ogni cella \(j=0,\dots,N-1\):
        \[
        f_j = \frac{1}{2} U_j^2.
        \]
        
        \item \textbf{Costante di Lax-Friedrichs:}
        \[
        \alpha = \max_{j} |U_j|.
        \]
        
        \item \textbf{Flux splitting:} per ogni cella \(j=0,\dots,N-1\):
        \[
        f^+_j = \frac{1}{2}\bigl(f_j + \alpha U_j\bigr), \qquad
        f^-_j = \frac{1}{2}\bigl(f_j - \alpha U_j\bigr).
        \]
        
        \item \textbf{Calcolo dei flussi numerici su tutti i bordi:}
        
        Per ogni bordo \(k = j+\tfrac12\) con \(j=0,\dots,N-1\):
        \begin{itemize}
            \item Si ricostruisce \(\hat f^+_{k}\) applicando la procedura WENO5 allo stencil sinistro (left-biased) sui valori:
            \[
            f^+_{j-2},\; f^+_{j-1},\; f^+_j,\; f^+_{j+1},\; f^+_{j+2}.
            \]
            \item Si ricostruisce \(\hat f^-_{k}\) applicando la procedura WENO5 allo stencil destro (right-biased) sui valori:
            \[
            f^-_{j-1},\; f^-_j,\; f^-_{j+1},\; f^-_{j+2},\; f^-_{j+3}.
            \]
            \item Il flusso numerico al bordo è la somma:
            \[
            F_{k} = \hat f^+_{k} + \hat f^-_{k}.
            \]
        \end{itemize}
        Valori con indici fuori dal dominio sono ottenuti dalle condizioni al contorno. \newline
        \textit{Poiché $f^+$ viaggia sempre a destra (ha derivata positiva), l'informazione per $f^+$ arriva da sinistra. Quindi per ricostruire $f^+$ al bordo $j+\tfrac12$ si usa uno stencil sbilanciato a sinistra. Analogamente, poiché $f^-$ viaggia sempre a sinistra (ha derivata negativa), l'informazione per $f^-$ arriva da destra; per ricostruire $f^-$ al bordo $j+\tfrac12$ si usa quindi uno stencil sbilanciato a destra.}
        \item \textbf{Calcolo di \(L(U)\):} per ogni cella \(j=0,\dots,N-1\):
        \[
        L(U)_j = -\frac{1}{\Delta x}\bigl(F_{j+1/2} - F_{j-1/2}\bigr).
        \]
    \end{enumerate}
    
    \item \textbf{Integrazione temporale SSP-RK3.}
    
    Dato lo stato corrente \(U^n\) al tempo \(t_n\), si eseguono i tre stadi:
    \[
    \begin{aligned}
    U^{(1)} &= U^n + \Delta t\, L(U^n), \\[4pt]
    U^{(2)} &= \frac34 U^n + \frac14 U^{(1)} + \frac14 \Delta t\, L(U^{(1)}), \\[4pt]
    U^{n+1} &= \frac13 U^n + \frac23 U^{(2)} + \frac23 \Delta t\, L(U^{(2)}).
    \end{aligned}
    \]
    Ogni chiamata di \(L(\cdot)\) ricalcola l'intera procedura al punto 1 sui valori correnti.
    
    \item \textbf{Condizioni al contorno.}
    
    Dopo ogni aggiornamento (sia nei sottostadi di Runge-Kutta sia al termine del passo temporale), si applicano le condizioni al contorno periodiche:
    \[
    U_{j+J} = U_j, \qquad \text{per ogni } j,
    \]
    dove \(J\) è il numero totale di celle del dominio.
\end{enumerate}

Il passo temporale \(\Delta t\) viene adattato dinamicamente secondo la condizione CFL:
\[
\Delta t = \text{CFL} \cdot \frac{\Delta x}{\max_j |U_j^n|},
\]
con \(\text{CFL} \le 0.5\) per garantire la stabilità di SSP-RK3 con lo splitting di Lax-Friedrichs.

Il metodo WENO5 con SSP-RK3 combina i pregi degli schemi precedenti: la conservatività esatta di Lax-Wendroff (che garantisce la corretta velocità dell'urto), l'assenza di oscillazioni tipica degli schemi monotoni come Godunov, e l'elevata accuratezza spaziale (quinto ordine) nelle regioni lisce, pur mantenendo un profilo d'urto estremamente localizzato su poche celle.

Riporto sotto una breve dimostrazione della necessità della condizione CFL:

La curva caratteristica \(X(t)\) soddisfa \(\frac{d}{dt}u(X(t),t)=0\). Sviluppando: \(u_t+u_x X'=0\). Confrontando con \(u_t+u u_x=0\) si ha \(X'=u\). Integrando \(dX = u\,dt\) tra \(t_0\) e \(t\):  
\(\int_{X(t_0)}^{X(t)} dX = \int_{t_0}^{t} u(X(s),s)\,ds\). Poiché \(u\) è costante lungo la caratteristica, \(u(X(s),s)=u(X(t_0),t_0)\), quindi  
\(X(t) = X(t_0) + u(X(t_0), t_0)(t-t_0)\).

Consideriamo il nodo \((X_j, t_{n+1})\). Scegliamo \(t_0 = t_{n+1}\) e \(X(t_{n+1})=X_j\). La caratteristica è  
\(X(t) = X_j + u(X_j, t_{n+1})(t-t_{n+1})\).  
Valutandola in \(t=t_n\):  
\(X(t_n) = X_j + u(X_j, t_{n+1})(t_n - t_{n+1})\).  
Poniamo \(x^* := X(t_n)\). Poiché \(u\) è costante lungo \(X(t)\), si ha \(u(X(t_n), t_n) = u(X(t_{n+1}), t_{n+1})\), cioè  
\(u(x^*, t_n) = u(X_j, t_{n+1})\).  
Sostituendo nella precedente e definendo \(\Delta t = t_{n+1}-t_n\):  
\(x^* = X_j - u(x^*, t_n)\,\Delta t\).  
Quindi il valore \(u(X_j, t_{n+1})\) dipende esclusivamente da \(u(x^*, t_n)\). Questo è il dominio di dipendenza fisico: un singolo punto \(x^*\) al tempo \(t_n\).

Ora passiamo allo schema numerico. Per calcolare \(u_j^{n+1}\), uno schema esplicito alle differenze finite usa un dominio di dipendenza numerico, cioè un insieme di nodi al tempo \(t_n\). Nel caso minimale a 3 punti (upwind o centrato), lo stencil è l'intervallo  
\([X_{j-1}, X_{j+1}]=[X_j-\Delta x, X_j+\Delta x]\).  
Affinché lo schema possa convergere alla soluzione esatta, il dominio fisico deve essere contenuto in quello numerico, altrimenti lo schema ignorerebbe l'informazione da cui la soluzione dipende realmente. Deve quindi valere  
\(X_j-\Delta x \le x^* \le X_j+\Delta x\).  
Sostituendo \(x^* = X_j - u(x^*, t_n)\Delta t\):  
\(-\Delta x \le -u(x^*, t_n)\Delta t \le \Delta x\),  
che equivale a  
\(|u(x^*, t_n)|\Delta t \le \Delta x\).  
Poiché la velocità \(u\) varia nello spazio e la disuguaglianza deve valere per ogni cella \(j\), si prende il caso peggiore sull'intero dominio:  
\(\max_j |u_j^n| \cdot \frac{\Delta t}{\Delta x} \le 1\).  
Da cui il vincolo  
\(\Delta t \le \frac{\Delta x}{\max_j |u_j^n|}\).  
Nella pratica si introduce un coefficiente di sicurezza \(\text{CFL} \le 0.5\) (dovuto alla stabilità effettiva di WENO5+SSP-RK3), ottenendo  
\(\Delta t = \text{CFL} \cdot \frac{\Delta x}{\max_j |u_j^n|}\).  
Questa condizione è necessaria: se violata, la caratteristica fisica esce dallo stencil numerico, e lo schema non può convergere alla soluzione debole corretta.
\newpage

\section{Soluzione esatta}

\subsection{Formulazione e problema del multivalore}

Come visto in Sezione~2.1, $u(x,t)=u_0(x_0)$ con $x_0$ soluzione di $x_0+u_0(x_0)t=x$. Dopo il tempo di rottura $T_b$ questa equazione può avere più soluzioni $x_0$ per lo stesso $x$: le caratteristiche si incrociano e la soluzione ``a più valori'' non ha senso fisico. Fra queste radici bisogna scegliere quella che realizza il minimo del funzionale di Lax--Oleinik (Sezione~2.4):
\[
J(y) = G(y) + \frac{(x-y)^2}{2t}, \qquad
u(x,t) = \frac{x-y^*}{t}, \qquad y^* = \arg\min_{y} J(y),
\]
dove $G(y)=\int^{y} u_0(s)\,ds$ è una primitiva di $u_0$ e $y$ varia su tutto $\mathbb{R}$ (non solo fra le radici). L'equivalenza con le caratteristiche è immediata:
\[
\frac{dJ}{dy} = u_0(y) - \frac{x-y}{t} = 0 \iff y+u_0(y)\,t = x,
\]
quindi i punti stazionari di $J$ sono le radici di $g(y)=y+u_0(y)t-x$; il minimo globale si trova confrontando i valori di $J$ in queste radici (e sui bordi/kink quando non ce ne sono, caso che descriveremo in Sezione~6.5).

\paragraph{Perché non ``il minimo di $u_0$''.} Un'euristica diffusa consiste nel prendere, fra le radici, quella con $u_0$ minimo. Con un solo urto isolato il risultato coincide con quello corretto, ma appena due urti interagiscono le radici provenienti da regioni lontane (dove $u_0$ è basso) possono avere il valore minimo pur non essendo la caratteristica che determina la soluzione fisica. Un controesempio numerico esplicito è in Sezione~6.4 e nel test \texttt{test\_soluzione\_esatta.py}: con $u_0=\sin(2\pi x)$ su $[-2,2]$ al tempo $t=0.5$, appena a sinistra dell'urto in $x=0.5$ la soluzione fisica vale $\approx 0.72$, mentre la regola del minimo di $u_0$ restituisce $\approx -0.77$.

Quanto segue descrive l'implementazione contenuta nel codice allegato (\texttt{exact\_solution} in \texttt{codice/burgers\_weno.py}), passo per passo.

\subsection{Caso $t\approx0$}

Se $t<10^{-14}$ la soluzione coincide banalmente con il dato iniziale: si valuta $u_0$ direttamente sui punti richiesti e si esce, evitando di applicare inutilmente (e con rischio di divisione per zero) la procedura di ricerca delle radici.

\subsection{Estensione del dominio di ricerca}

Per trovare tutte le radici $x_0$ dell'equazione implicita occorre esplorare un intervallo di $x_0$ sufficientemente ampio: le caratteristiche possono trasportare informazione da punti anche molto lontani dal dominio $[x_{\min},x_{\max}]$ di interesse. Si campiona $u_0$ su un intervallo esteso di lunghezza $3L$ (con $L=x_{\max}-x_{\min}$) per stimarne la velocità massima $\texttt{maxU0}$, e si pone
\[
\text{ext} = \max(\text{ext\_factor}\cdot L,\ 2\,\texttt{maxU0}\cdot t) + 2.0,
\]
dove il primo termine garantisce un'estensione minima pari a $3L$, il secondo copre la distanza massima percorribile dalle caratteristiche nel tempo $t$, e il margine $+2.0$ assorbe eventuali imprecisioni di campionamento.

\subsection{Ricerca delle radici e selezione entropica}

Definita $g(x_0)=x_0+u_0(x_0)t-x$, si valuta $g$ su una griglia fine di $N_s=2000$ punti $x_s$ nel dominio esteso, si individuano i cambi di segno tra punti adiacenti (candidati zero) e si raffina ciascuno con 64 iterazioni di \textbf{bisezione}, ottenendo una precisione ben oltre quella necessaria. Per ogni radice $y$ trovata si valuta il funzionale entropico
\[
J(y) = G(y) + \frac{(x-y)^2}{2t}
\]
(la primitiva $G$ è calcolata una volta per tutte sulla griglia estesa con la regola del trapezio e interpolata linearmente), e si tiene la radice con $J$ minimo:
\[
y^* = \arg\min\{J(y) : g(y)=0\}, \qquad u(x,t) = \frac{x-y^*}{t}.
\]
L'uso di $J$ al posto del valore $u_0(y)$ è l'unica differenza rispetto alla versione precedente della soluzione esatta: nella prima versione si prendeva la radice con $u_0$ minimo, che non è in generale lo stato entropico.
La funzione gestisce inoltre eventuali \texttt{NaN} del dato iniziale sostituendoli con 0, garantendo robustezza anche per condizioni iniziali complesse come quelle usate in questo progetto.

\paragraph{Controesempio numerico (perché la regola del minimo di $u_0$ è sbagliata).} Si consideri $u_0(x)=\sin(2\pi x)$ su $[-2,2]$, che al tempo $t=0.5$ ha un urto in $x=0.5$ (per simmetria). Subito a sinistra dell'urto, $x=0.486$, l'equazione $g(y)=0$ ha più radici: quella fisica fornisce
\[
u(x,t)=\frac{x-y^*}{t}\approx 0.72,
\]
mentre la radice con $u_0$ minimo (una caratteristica partita da molto lontano) dà $\approx -0.77$. Un riferimento numerico ad alta risoluzione (schema conservativo indipendente, $2000$ celle) conferma il primo valore: è il test \texttt{test\_stato\_entropico\_accanto\_all\_urto} di \texttt{test\_soluzione\_esatta.py}.

\paragraph{Esempio con l'IC del progetto.} Per $u_0(x)=-\sin(2x)$ su $[-\pi,\pi]$ il tempo di rottura è $T_b=\frac{1}{\max(-u_0')}=\frac12$: per $t=0.4<T_b$ la soluzione è ancora a un solo valore e le due regole coincidono (massima differenza misurata $<10^{-6}$), mentre per $t=1$ gli urti si sono formati e la differenza rispetto al riferimento indipendente è di due ordini di grandezza maggiore per la vecchia regola (si vedano i test e la Sezione~2.4).
\newpage
\subsection{Caso particolare: rarefazione di Riemann}

L'algoritmo di ricerca delle radici descritto in Sezione~6.4 funziona correttamente per tutti i dati iniziali continui e per gli urti. Tuttavia, se la condizione iniziale presenta una \emph{discontinuità verso l'alto} (rarefazione), l'equazione \(x_0+u_0(x_0)t=x\) \textbf{non ha soluzioni} nella regione del ventaglio: nessuna caratteristica arriva in quei punti. In questo caso il funzionale $J$ non ha punti stazionari e il suo minimo cade sul \emph{kink} della primitiva $G$ (in corrispondenza del salto della condizione iniziale): l'implementazione descritta in Sezione~6.4 usa allora il punto di minimo di $J$ sulla griglia e restituisce proprio il valore del ventaglio, $u=x/t$. Conviene comunque ricavare esplicitamente la soluzione del ventaglio, considerando il problema di Riemann puro:

\[
u_t + u u_x = 0, \qquad 
u(x,0)=
\begin{cases}
u_l, & x<0,\\
u_r, & x>0,
\end{cases}
\qquad \text{con } u_l<u_r.
\]

Poiché la condizione iniziale è costante a tratti, il problema è invariante per scalatura: se moltiplichiamo \(x\) e \(t\) per una stessa costante \(\lambda\), l'equazione e i dati iniziali restano identici. La soluzione deve quindi dipendere solo dal rapporto \(\xi=x/t\). Poniamo \(u(x,t)=\phi(\xi)\). Allora:

\[
u_t = -\frac{\xi}{t}\phi', \qquad u_x = \frac{1}{t}\phi',
\]
e sostituendo in \(u_t+u u_x=0\):
\[
-\frac{\xi}{t}\phi' + \phi\frac{1}{t}\phi' = 0
\;\Longrightarrow\; \phi'(\phi-\xi)=0.
\]

Le due soluzioni elementari sono:
\[
\phi' = 0 \;\Longrightarrow\; \phi=\text{costante}, \qquad \text{oppure} \qquad \phi=\xi.
\]
Per raccordare le zone piatte \(\phi=u_l\) (a sinistra) e \(\phi=u_r\) (a destra) senza creare un salto (condizione di entropia di Oleinik per flussi convessi), si deve usare il ramo \(\phi=\xi\) nel mezzo. La soluzione debole continua è quindi:

\[
u(x,t) = 
\begin{cases}
u_l, & \dfrac{x}{t} \le u_l, \\[6pt]
\dfrac{x}{t}, & u_l \le \dfrac{x}{t} \le u_r, \\[6pt]
u_r, & \dfrac{x}{t} \ge u_r.
\end{cases}
\tag{*}
\]

Questa è l'onda di rarefazione centrata: un ventaglio di caratteristiche che partono tutte dall'origine \(x=0\) al tempo \(t=0\), con velocità \(s\in[u_l,u_r]\). Notiamo che la formula \((*)\) è esattamente il caso limite della formula di Lax--Oleinik: se il salto iniziale è in \(y_0\), la primitiva \(G\) ha il kink in \(y_0\) e il minimo di \(J\) cade lì, dando \(u=(x-y_0)/t=x/t\) (per \(y_0=0\)). Per dati iniziali generici (non costanti a tratti) il ventaglio si ottiene con la stessa formula, senza bisogno di alcun trattamento speciale: il codice allegato la calcola automaticamente, come verificato dal test \texttt{test\_ventaglio\_di\_rarefazione}.
\newpage
\textit{Sappiamo che per un punto $x(t)$ fissato abbiamo un valore $u$; sappiamo che da quel punto $x(t)$ esiste una retta nel tempo per cui tutti i valori $u$ sono uguali. Quindi dalla condizione iniziale possiamo trovare esattamente i valori della soluzione esatta per ogni $t$ e per ogni $x$ semplicemente espandendo tantissime rette. Per farlo vediamo che $x(t)=x_0+u_0(x_0)t$; a questa equazione si arriva integrando $dx/dt = u$ e sfruttando il fatto che $u$ deve essere costante lungo $dx/dt$, quindi $dx/dt=u_0$; poi si integra $x$ da un tempo $t_0$ a un tempo $t$ e si sposta $x_0$ dall'altra parte. Una volta che abbiamo quell'equazione della retta $x(t)$, possiamo trovare un qualsiasi valore di $u(x(t),t)$ solamente sapendo $x_0$, ovvero il punto $x$ nella soluzione iniziale $u_0$ da cui si propaga quella caratteristica, cioè la partenza della retta. Una volta trovato quel punto, basta andare a vedere che valore avevamo in $u_0(x_0)$ e copiarlo al punto $(x(t),t)$. In corrispondenza di un urto si hanno più valori di $x_0$ da cui partono le caratteristiche: scegliamo semplicemente il minimo di questi, e questo punto rappresenta la soluzione fisica (soluzione dell'equazione con viscosità ma con $v$ che tende a zero). In corrispondenza di una rarefazione (problema di Riemann con discontinuità nelle condizioni iniziali) NON si ha un punto $x_0$ tale per cui la caratteristica arriva al punto $(x(t),t)$; allora semplicemente poniamo $u(x,t)=x/t$.}
\newpage

\begin{comment}

\section{Implementazione e verifica numerica}

\subsection{Struttura del codice}

Il codice Python implementa, oltre alla soluzione esatta descritta in Sezione~6, il solutore WENO generico di ordine $2k-1$ per $k=1,2,3$ (rispettivamente upwind, WENO3, WENO5), integrato in tempo con SSP-RK3 (Sezione~5.2) e con passo temporale adattato dinamicamente alla condizione CFL,
\[
\Delta t = \text{CFL}\cdot\frac{\Delta x}{\max_j|U_j^n|}.
\]
La condizione iniziale usata nel progetto è volutamente più complessa del semplice $\sin(\pi x)$: $u_0(x)=\sin(\pi x^6-4x^3+2x)$, cosicché la soluzione esatta sviluppi più urti e regioni lisce di curvatura diversa, mettendo alla prova sia la ricerca delle radici (Sezione~6.4) sia la capacità di WENO di adattarsi a geometrie non banali.

\subsection{Test di convergenza}

Il codice confronta, a parità di griglia e tempo finale, l'errore in norma $L^1$ tra soluzione esatta e soluzione numerica per $k=1,2,3$:
\[
\|e_k\|_{L^1} = \Delta x\sum_j |u_{\text{esatta}}(x_j)-U_j^{k}|,
\]
e ne calcola il rapporto tra schemi successivi. In una zona con urti attivi non ci si aspetta di osservare l'ordine asintotico pieno (5 per WENO5), perché la presenza della discontinuità limita la convergenza a ordine 1 localmente; ci si aspetta però una gerarchia monotona degli errori, $\|e_3\|<\|e_2\|<\|e_1\|$, che conferma che l'aumento dell'ordine di ricostruzione migliora comunque la qualità globale della soluzione (minore diffusione numerica lontano dagli urti, urto più netto). Il codice stampa un verdetto automatico (\texttt{PASS}, \texttt{PARZIALE} o \texttt{SATURATO}) basato su questa gerarchia.

\subsection{Interpretazione dei risultati: forma attesa dell'urto numerico}

Nell'analisi dei grafici è cruciale distinguere tra errore di posizione ed errore di forma. Su una griglia finita, nessuno schema conservativo può riprodurre un salto verticale perfetto, perché la soluzione numerica è definita solo su un insieme discreto di punti e la ricostruzione polinomiale non può rappresentare una discontinuità senza generare oscillazioni (fenomeno di Runge). WENO5 risolve questo compromesso producendo, in corrispondenza dell'urto, una transizione molto ripida ma comunque spalmata su 2–3 celle, anziché un gradino netto. 

Questa "rampa" non è un errore di posizione: grazie alla struttura conservativa dello schema, la massa totale e la velocità dell'urto sono quelle corrette (Rankine-Hugoniot), e il salto numerico cade esattamente nel punto previsto dalla soluzione esatta. La differenza rispetto alla soluzione esatta è quindi puramente locale (sull'urto) e si manifesta in un errore puntuale \(O(1)\) in quel ristretto intorno. Tuttavia, l'area sottesa alla curva numerica coincide con quella esatta; per questo motivo l'errore globale in norma \(L^1\) scala come \(O(\Delta x)\) (ordine 1), nonostante la ricostruzione spaziale sia di quinto ordine nelle zone lisce. Al raffinare della griglia, la larghezza di questa rampa si riduce proporzionalmente a \(\Delta x\), e la soluzione numerica converge uniformemente alla discontinuità nel senso delle distribuzioni. 

Rispetto agli schemi di ordine inferiore, WENO5 offre quindi il miglior compromesso possibile: un urto netto quanto la griglia lo consente (molto più di Upwind), privo di oscillazioni (a differenza di Lax-Wendroff) e correttamente posizionato (grazie alla conservatività). Questo comportamento, apparentemente "meno netto" della soluzione esatta, è la migliore approssimazione raggiungibile su una risoluzione spaziale finita senza violare i vincoli di stabilità ed entropia.
\section{Conclusioni}

Il progetto mostra, in modo autocontenuto, l'intera catena logica che porta dagli schemi classici a un metodo moderno per leggi di conservazione: (i) l'equazione di Burgers sviluppa urti anche da dati lisci; (ii) in presenza di urti la nozione classica di soluzione va sostituita da quella di soluzione debole entropica, caratterizzata da Rankine-Hugoniot e dalla condizione di Oleinik; (iii) uno schema costruito applicando Taylor punto per punto (anche se stabile) non conserva la massa discreta e produce una velocità d'urto sbagliata; (iv) la forma conservativa risolve il problema per costruzione, tramite la cancellazione telescopica dei flussi; (v) il metodo WENO5 realizza la forma conservativa con un flusso numerico che è di quinto ordine nelle zone lisce e degrada automaticamente vicino agli urti, eliminando le oscillazioni di Gibbs senza sacrificare l'accuratezza lontano da essi. Il confronto numerico con la soluzione esatta, calcolata indipendentemente tramite il metodo delle caratteristiche e la condizione di Oleinik, conferma questa costruzione teorica.
\end{comment}
\newpage
\begin{pythoncode}[title=Implementazione WENO]

#!/usr/bin/env python3
# -*- coding: utf-8 -*-
"""Burgers 1D — schema WENO5 + SSP-RK3 e soluzione esatta entropica.

Contiene tutto quello che serve al progetto di Calcolo Scientifico
(Progetto 14) "High order FD discretization of Burgers' equation in 1D with
WENO":

  * il solutore numerico di  u_t + (u^2/2)_x = 0  con
      - flux splitting di Lax-Friedrichs,
      - ricostruzione WENO5 alle interfacce di cella,
      - avanzamento temporale SSP-RK3 con dt adattato alla CFL;
  * la soluzione esatta entropica, calcolata con la formula di
    **Hopf-Lax-Oleinik**:
        u(x,t) = (f*)'(p*),   p* = (x - y*)/t,
        y* = argmin_y [ G(y) + t f*((x-y)/t) ],   G(y) = int u0
    che per Burgers (f*(p) = p^2/2) diventa
        y* = argmin_y [ G(y) + (x-y)^2/(2t) ],    u(x,t) = (x - y*)/t.
    I punti stazionari di J(y) sono le radici dell'equazione delle
    caratteristiche y + u0(y) t = x; fra essi si sceglie pero' quello che
    MINIMIZZA J (condizione di entropia), non quello con u0 minimo.

Uso:
    python3 burgers_weno.py                 # errori L1 + grafico su file
    python3 burgers_weno.py --show          # come sopra, ma apre la finestra
    python3 burgers_weno.py --nx 800 --t 0.9

La funzione `old_min_rule` (in fondo) riproduce la versione precedente —
`u = min{ u0(y) : y + u0(y) t = x }` — e serve solo a documentare, su casi con
piu' urti che interagiscono, che quella regola NON seleziona la soluzione
entropica (si vedano README e test_soluzione_esatta.py).
"""

import argparse
import os

import numpy as np

# ----------------------------------------------------------------------------
# Parametri di default (quelli del progetto)
# ----------------------------------------------------------------------------
AMP = 1.0          # ampiezza della condizione iniziale
XMIN, XMAX = -np.pi, np.pi
NX = 200           # celle della griglia numerica
CFL = 0.45         # coefficiente di Courant (limite pratico 0.5 per WENO5+RK3)
T_FINAL = 0.4      # tempo finale della simulazione
NX_EXACT = 2000    # punti della griglia di ricerca per la soluzione esatta
EXT_FACTOR = 3.0   # estensione minima del dominio di ricerca (in unita' di L)

PAD = 3            # ghost cells necessarie allo stencil di 5 punti di WENO5
EPS = 1e-36        # stabilizzatore dei pesi non lineari


def u0(x):
    """u(x,0) = -sin(2x): la IC del progetto (piu' urti che interagiscono e
    regioni di curvatura diversa)."""
    return AMP * np.sin(np.pi + 2 * x)


# ----------------------------------------------------------------------------
# WENO5 + SSP-RK3
# ----------------------------------------------------------------------------
def weno_reconstruct(f, n, side):
    """Ricostruzione WENO5 del flusso all'interfaccia j+1/2.

    side='m' usa lo stencil orientato a sinistra (f_{j-2}...f_{j+2}),
    side='p' lo stencil orientato a destra   (f_{j-1}...f_{j+3});
    i pesi lineari ottimali si scambiano nei due casi.
    """
    fpad = np.pad(f, (PAD, PAD), mode='wrap')
    if side == 'm':
        v0 = fpad[PAD - 2: PAD - 2 + n]   # f_{j-2}
        v1 = fpad[PAD - 1: PAD - 1 + n]   # f_{j-1}
        v2 = fpad[PAD    : PAD     + n]   # f_j
        v3 = fpad[PAD + 1: PAD + 1 + n]   # f_{j+1}
        v4 = fpad[PAD + 2: PAD + 2 + n]   # f_{j+2}
        p0 = (1.0 / 3.0) * v0 - (7.0 / 6.0) * v1 + (11.0 / 6.0) * v2
        p1 = -(1.0 / 6.0) * v1 + (5.0 / 6.0) * v2 + (1.0 / 3.0) * v3
        p2 = (1.0 / 3.0) * v2 + (5.0 / 6.0) * v3 - (1.0 / 6.0) * v4
        b0 = (13.0 / 12.0) * (v0 - 2 * v1 + v2) ** 2 + (1.0 / 4.0) * (v0 - 4 * v1 + 3 * v2) ** 2
        b1 = (13.0 / 12.0) * (v1 - 2 * v2 + v3) ** 2 + (1.0 / 4.0) * (v1 - v3) ** 2
        b2 = (13.0 / 12.0) * (v2 - 2 * v3 + v4) ** 2 + (1.0 / 4.0) * (3 * v2 - 4 * v3 + v4) ** 2
        d0, d1, d2 = 0.1, 0.6, 0.3
    else:                                  # side == 'p'
        v0 = fpad[PAD - 1: PAD - 1 + n]   # f_{j-1}
        v1 = fpad[PAD    : PAD     + n]   # f_j
        v2 = fpad[PAD + 1: PAD + 1 + n]   # f_{j+1}
        v3 = fpad[PAD + 2: PAD + 2 + n]   # f_{j+2}
        v4 = fpad[PAD + 3: PAD + 3 + n]   # f_{j+3}
        p0 = (1.0 / 3.0) * v0 - (7.0 / 6.0) * v1 + (11.0 / 6.0) * v2
        p1 = -(1.0 / 6.0) * v1 + (5.0 / 6.0) * v2 + (1.0 / 3.0) * v3
        p2 = (1.0 / 3.0) * v2 + (5.0 / 6.0) * v3 - (1.0 / 6.0) * v4
        b0 = (13.0 / 12.0) * (v0 - 2 * v1 + v2) ** 2 + (1.0 / 4.0) * (v0 - 4 * v1 + 3 * v2) ** 2
        b1 = (13.0 / 12.0) * (v1 - 2 * v2 + v3) ** 2 + (1.0 / 4.0) * (v1 - v3) ** 2
        b2 = (13.0 / 12.0) * (v2 - 2 * v3 + v4) ** 2 + (1.0 / 4.0) * (3 * v2 - 4 * v3 + v4) ** 2
        d0, d1, d2 = 0.3, 0.6, 0.1
    w0 = d0 / (b0 + EPS) ** 2
    w1 = d1 / (b1 + EPS) ** 2
    w2 = d2 / (b2 + EPS) ** 2
    return (w0 * p0 + w1 * p1 + w2 * p2) / (w0 + w1 + w2)


def compute_L(u, dx):
    """Residuo spaziale L(u) = -(F_{j+1/2} - F_{j-1/2})/dx (forma conservativa)."""
    n = len(u)
    f = 0.5 * u * u                       # flusso fisico f(u) = u^2/2
    alpha = np.max(np.abs(u))             # velocita' caratteristica massima
    fp = 0.5 * (f + alpha * u)            # splitting di Lax-Friedrichs
    fm = 0.5 * (f - alpha * u)
    F_hat = weno_reconstruct(fp, n, 'm') + weno_reconstruct(fm, n, 'p')
    F_ext = np.concatenate([[F_hat[-1]], F_hat])   # condizioni periodiche
    return -(F_ext[1:] - F_ext[:-1]) / dx


def ssp_rk3(u, dx, dt):
    """Terzo ordine SSP (Shu-Osher) in tre stadi."""
    L0 = compute_L(u, dx)
    u1 = u + dt * L0
    L1 = compute_L(u1, dx)
    u2 = 0.75 * u + 0.25 * (u1 + dt * L1)
    L2 = compute_L(u2, dx)
    return (1.0 / 3.0) * u + (2.0 / 3.0) * (u2 + dt * L2)


def solve_weno(u0_arr, dx, T_final, cfl):
    """Integra in tempo con dt adattato alla condizione CFL."""
    u = u0_arr.copy()
    t = 0.0
    while t < T_final:
        maxU = np.max(np.abs(u))
        dt = cfl * dx / maxU if maxU > 1e-14 else 1e-12
        if t + dt > T_final:
            dt = T_final - t
        u = ssp_rk3(u, dx, dt)
        t += dt
    return u


# ----------------------------------------------------------------------------
# SOLUZIONE ESATTA — formula di Hopf-Lax-Oleinik
# ----------------------------------------------------------------------------
def primitiva(u0_func, ys, u0sy=None):
    """G(y) = int_{ys[0]}^{y} u0(s) ds, con la regola del trapezio."""
    if u0sy is None:
        u0sy = np.nan_to_num(np.array([u0_func(y) for y in ys]), nan=0.0)
    G = np.zeros(np.size(ys))
    G[1:] = np.cumsum(0.5 * (u0sy[1:] + u0sy[:-1]) * np.diff(ys))
    return G


def exact_solution(u0_func, x_vec, t, Ns=NX_EXACT, ext_factor=EXT_FACTOR):
    """Soluzione entropica di Burgers (formula di Hopf-Lax-Oleinik).

    Per f(u) = u^2/2 si ha f*(p) = p^2/2, quindi

        u(x,t) = (x - y*)/t,   y* = argmin_y [ G(y) + (x-y)^2 / (2t) ].

    I punti stazionari del funzionale J(y) = G(y) + (x-y)^2/(2t) sono le
    radici di  g(y) = y + u0(y) t - x  (le caratteristiche che arrivano in x);
    fra i candidati si sceglie quello che MINIMIZZA J — questa e' la
    condizione di entropia, e non equivale a prendere il minimo di u0.
    Se nessuna caratteristica arriva in x (ventaglio di rarefazione) il minimo
    di J cade sul bordo/kink della primitiva e la stessa formula restituisce
    u = x/t, cioe' la nota soluzione del ventaglio.
    """
    x_vec = np.asarray(x_vec, dtype=float)
    if t < 1e-14:                                    # t = 0: la IC stessa
        return np.array([u0_func(xi) for xi in x_vec])

    xmin, xmax = x_vec[0], x_vec[-1]
    L = xmax - xmin
    # max|u0| = max|f'(u0)| sul dominio esteso -> dimensione dell'estensione
    maxU0 = 0.1
    Nc = 500
    for j in range(Nc + 1):
        v = u0_func(xmin - L + j * (3.0 * L / Nc))
        if np.isfinite(v) and abs(v) > maxU0:
            maxU0 = abs(v)
    ext = max(ext_factor * L, 2.0 * maxU0 * t) + 2.0
    xs = np.linspace(xmin - ext, xmax + ext, Ns)
    u0s = np.nan_to_num(np.array([u0_func(y) for y in xs]), nan=0.0)
    G = primitiva(u0_func, xs, u0s)

    def G_at(y):
        i = np.clip(np.searchsorted(xs, y), 1, Ns - 1)
        frac = (y - xs[i - 1]) / (xs[i] - xs[i - 1])
        return G[i - 1] + frac * (G[i] - G[i - 1])

    def J(y, x):
        return G_at(y) + (x - y) ** 2 / (2.0 * t)

    out = np.zeros(x_vec.size)
    for idx, x in enumerate(x_vec):
        gs = xs + u0s * t - x
        # (1) candidati: zeri di g(y) = 0 (punti stazionari di J, bisezione)
        best_y, best_J = None, np.inf
        for j in np.where(gs[:-1] * gs[1:] <= 0)[0]:
            lo, hi, glo = xs[j], xs[j + 1], gs[j]
            for _ in range(64):
                mid = 0.5 * (lo + hi)
                gm = mid + u0_func(mid) * t - x
                if glo * gm <= 0:
                    hi = mid
                else:
                    lo, glo = mid, gm
            y = 0.5 * (lo + hi)
            Jv = J(y, x)
            if Jv < best_J:
                best_y, best_J = y, Jv
        # (2) fallback: se nessuna caratteristica arriva in x (vento di
        #     rarefazione) si minimizza J sulla griglia, con raffinamento
        if best_y is None:
            j0 = int(np.argmin(G + (x - xs) ** 2 / (2.0 * t)))
            a, b = xs[max(0, j0 - 1)], xs[min(Ns - 1, j0 + 1)]
            phi = (np.sqrt(5.0) - 1.0) / 2.0
            c, d = b - phi * (b - a), a + phi * (b - a)
            fc, fd = J(c, x), J(d, x)
            for _ in range(80):
                if fc < fd:
                    b, d, fd = d, c, fc
                    c = b - phi * (b - a)
                    fc = J(c, x)
                else:
                    a, c, fc = c, d, fd
                    d = a + phi * (b - a)
                    fd = J(d, x)
            best_y = 0.5 * (a + b)
        out[idx] = (x - best_y) / t                  # u(x,t) = (x - y*)/t
    return out


def old_min_rule(u0_func, x_vec, t, Ns=NX_EXACT, ext_factor=EXT_FACTOR):
    """Versione PRECEDENTE (non corretta) della soluzione esatta:

        u(x,t) = min{ u0(y) : y + u0(y) t = x }.

    Serve solo a documentare l'errore corretto in questo progetto: quando piu'
    urti interagiscono, il minimo di u0 fra le caratteristiche che arrivano in
    (x,t) non e' la selezione entropica (si veda test_soluzione_esatta.py).
    """
    x_vec = np.asarray(x_vec, dtype=float)
    if t < 1e-14:
        return np.array([u0_func(xi) for xi in x_vec])
    xmin, xmax = x_vec[0], x_vec[-1]
    L = xmax - xmin
    maxU0 = 0.1
    Nc = 500
    for j in range(Nc + 1):
        v = u0_func(xmin - L + j * (3.0 * L / Nc))
        if np.isfinite(v) and abs(v) > maxU0:
            maxU0 = abs(v)
    ext = max(ext_factor * L, 2.0 * maxU0 * t) + 2.0
    xs = np.linspace(xmin - ext, xmax + ext, Ns)
    u0s = np.nan_to_num(np.array([u0_func(y) for y in xs]), nan=0.0)
    out = np.zeros(x_vec.size)
    for idx, x in enumerate(x_vec):
        gs = xs + u0s * t - x
        cands = []
        for j in np.where(gs[:-1] * gs[1:] <= 0)[0]:
            lo, hi, glo = xs[j], xs[j + 1], gs[j]
            for _ in range(64):
                mid = 0.5 * (lo + hi)
                gm = mid + u0_func(mid) * t - x
                if glo * gm <= 0:
                    hi = mid
                else:
                    lo, glo = mid, gm
            u0c = u0_func(0.5 * (lo + hi))
            if np.isfinite(u0c):
                cands.append(u0c)
        out[idx] = min(cands) if cands else u0s[int(np.argmin(np.abs(gs)))]
    return out


# ----------------------------------------------------------------------------
# Main
# ----------------------------------------------------------------------------
def main():
    ap = argparse.ArgumentParser(
        description="Burgers 1D: WENO5 + SSP-RK3 contro la soluzione esatta di Hopf-Lax-Oleinik")
    ap.add_argument("--nx", type=int, default=NX, help="celle della griglia numerica")
    ap.add_argument("--cfl", type=float, default=CFL, help="coefficiente CFL")
    ap.add_argument("--t", type=float, default=T_FINAL, help="tempo finale")
    ap.add_argument("--nx-exact", type=int, default=NX_EXACT,
                    help="punti della griglia di ricerca della soluzione esatta")
    ap.add_argument("--show", action="store_true", help="apre la finestra del grafico")
    ap.add_argument("--no-figure", action="store_true", help="non salva il grafico")
    args = ap.parse_args()

    import matplotlib
    if not args.show:
        matplotlib.use("Agg")
    import matplotlib.pyplot as plt

    x, dx = np.linspace(XMIN, XMAX, args.nx, endpoint=False, retstep=True)
    x = x + 0.5 * dx
    u_ini = u0(x)
    u_exact = exact_solution(u0, x, args.t, args.nx_exact, EXT_FACTOR)
    u_weno = solve_weno(u_ini, dx, args.t, args.cfl)
    diff = np.abs(u_weno - u_exact)
    err_medio = float(diff.mean())
    err_l1_rel = float(diff.sum() / np.abs(u_exact).sum())

    print(f"Nx={args.nx}  CFL={args.cfl}  T={args.t}")
    print(f"errore medio |u_num - u_esatta|  = {err_medio:.6e}")
    print(f"errore L1 relativo (vs ||u_esatta||_1) = {err_l1_rel:.6e}")

    plt.figure(figsize=(10, 6))
    plt.plot(x, u_ini, 'k--', linewidth=1.5, label='t=0')
    plt.plot(x, u_exact, 'b-', linewidth=2.5, label='Esatta (Hopf-Lax-Oleinik)')
    plt.plot(x, u_weno, 'r-', linewidth=2, label='WENO5')
    plt.xlabel('x')
    plt.ylabel('u(x,t)')
    plt.title(f'Burgers, t={args.t:.3f}, Nx={args.nx}, CFL={args.cfl}')
    plt.legend()
    plt.grid(alpha=0.3)
    if not args.no_figure:
        out = os.path.normpath(os.path.join(os.path.dirname(os.path.abspath(__file__)),
                                            os.pardir, "figure"))
        os.makedirs(out, exist_ok=True)
        path = os.path.join(out, "burgers_weno.png")
        plt.savefig(path, dpi=130, bbox_inches="tight")
        print(f"grafico salvato in {path}")
    if args.show:
        plt.show()


if __name__ == "__main__":
    main()
\end{pythoncode}
\end{document}

